\section{Conclusion and Future Work}
\label{sec:conclusion}

Without a memory image, three parts of migration theory still apply, the split into cold and live migration, the downtime and the total migration time. However, the strategies have to transfer files instead of memory pages. All four migration types of He and Buyya have a counterpart at the filesystem level, each with its own condition (RQ-1). Pre-copy requires an anticipated handover, and post-copy needs the browser tab to stay open until the pull phase ends. Hybrid migration depends on both of these conditions. Pre-copy and cold migration are the two types measured in this work. In the live study at the large level, both push strategies create the target before the stop and use a single push round. This lowers the median interruption time of the heartbeat from \SI{8.74}{\second} to \SI{1.69}{\second} for pre-copy and from \SI{6.91}{\second} to \SI{1.57}{\second} for manifest-push. However, pre-copy takes longer than the full copy in total. Without anticipation, manifest-based transfer beats the full copy from a project size between the medium and the large level (RQ-2). For these four projects on E2B, this size lies between \SI{4.62}{\mega\byte} and \SI{13.93}{\mega\byte}. It may also depend on their packages, on the uplink of the client and on the registry link of the target.

In future work, the handover could be compared with a cooperative case, in which the runtime can be controlled and the memory can come along. With such a comparison, the cost of the missing memory image could be measured. How much the choice of source and target affects the results could be measured with other sources and targets. Furthermore, post-copy and hybrid migration could be measured after the first tick with a workload that reads its files. Additionally, it could be measured whether and how often a session reaches the limits of the sandbox in practice and whether the source can observe this in advance.
